% Propietario conservado: /Users/ruben/Documents/New project/output/EXTREMA_MEDIA_RAZON_RECIPROCIDAD_ES_EN_20260918/ARTICULO/ES/source/libro/06_sintesis.tex
\chapter{Unidad dimensional, simetría relativa y recuperación de la estructura}
\label{x_reciprocidad:lib06:sintesis}

La extrema y media razón holográfica articula la conservación evolutiva
de la unidad con sus realizaciones aritméticas y dimensionales. La
composición desarrollada permite precisar ese principio mediante
cantidades conservadas, dominios e inversas. El estado enriquecido
generado por APP--TRIT--TPK precede a las lecturas numéricas y
geométricas; sus operaciones conservan la orientación, el acarreo,
la incidencia y la memoria que cada realización utiliza.

\section{Conservación aritmética, unitariedad y descomposición isométrica}
\label{x_reciprocidad:lib06:tres-conservaciones}

El refinamiento aritmético conserva la unidad normalizada:
\begin{equation}
 (x,y,c)\longmapsto(Bx-c,By+c),\qquad
 \frac{Bx-c}{B(x+y)}+\frac{By+c}{B(x+y)}=1.
 \label{x_reciprocidad:lib06:unidad}
\end{equation}
La inversa por residuos recupera el acarreo canónico y la pareja
anterior. La composición de bloques conserva las contribuciones de
los eventos con su escala, y las divisiones sucesivas recuperan toda
historia finita. La lectura límite se controla mediante cilindros
decrecientes; las marcas de frontera mantienen la distinción entre
historias que comparten una evaluación escalar.

El levantamiento sobre etiquetas transforma esa biyección en una
unitaria \(J\). Conserva la norma de las amplitudes y transporta
los lectores de fracción y acarreo mediante
\(J^*M_hJ=M_{h\circ\ell}\). La cantidad \(x^2+y^2\) y la norma
de \(|x,y,c\rangle\) pertenecen a realizaciones distintas. La
prueba de unitariedad establece el enlace preciso entre la
reversibilidad aritmética y la conservación de amplitudes.

La tercera conservación corresponde a dos realizaciones isométricas:
\begin{equation}
 (a+1)Q_-^*Q_-+aQ_+^*Q_+
 =(2a+1)(a^2+a+1)I.
 \label{x_reciprocidad:lib06:balance}
\end{equation}
La cancelación de los términos mixtos permite la distribución
bilateral y su reconstrucción de norma uno. La ley vale para los
transportes del dominio especificado en la
\cref{x_reciprocidad:lib04:teo-balance}. Su especialización al ciclo ternario
aporta propiedades adicionales de norma e índice; la ley general
conserva su validez para isometrías que carezcan de ese orden finito.

La cadena formada por las tres aplicaciones es efectiva: la biyección
produce \(J\); una permutación previamente declarada produce el
segundo transporte \(JR\); los dos transportes satisfacen el
balance, y sus canales recuperan la entrada por
\eqref{x_reciprocidad:lib05:inversa-bilateral}. El ejemplo de la
\cref{x_reciprocidad:lib05:ejemplo} muestra el procedimiento completo con
\((73,27,1)\), \((729,271)\) y \(1001\) etiquetas admisibles.
Los números de esa realización no se introducen como sustitutos de
la estructura que genera el registro \(K\).

\section{Memoria completa y contribución lógica al límite de Landauer}
\label{x_reciprocidad:landauer:puente-archivo}

La contribución lógica nula de Landauer, establecida en
\eqref{x_reciprocidad:v:ant:eq:ii-kb-aut-reversible}, se aplica a los transportes
recuperables construidos en este artículo. La biyección de
\cref{x_reciprocidad:lib05:teo-biyeccion} conserva el acarreo por
\eqref{x_reciprocidad:lib05:inversa-etiquetas}; su prolongación conserva la ruta
mediante \cref{x_reciprocidad:lib05:teo-rutas}. Los archivos de
\cref{x_reciprocidad:lib04:teo-archivo} conservan las amplitudes y correlaciones
necesarias para invertir el transporte.

\begin{corollary}[Recuperación y conservación espectral del archivo]
\label{x_reciprocidad:landauer:cor-archivo}
Sea \(Z:H\to\mathcal K\) uno de los archivos completos
\(Z_N,Z_\infty\) de \cref{x_reciprocidad:lib04:teo-archivo}, o el levantamiento
unitario \(J_n\) de \cref{x_reciprocidad:lib05:teo-rutas}, en su respectivo dominio.
Para un operador de densidad de rango finito \(\rho\) sobre \(H\),
la salida \(\rho_Z=Z\rho Z^*\) satisface
\begin{align}
 \operatorname{Tr}\rho_Z&=1,& Z^*\rho_ZZ&=\rho,\\
 s(\rho_Z)&=s(\rho),&
 s(\rho)&=-\operatorname{Tr}(\rho\log\rho),
 \label{x_reciprocidad:landauer:espectro}
\end{align}
con \(0\log0=0\). Si \(X\) es una variable aleatoria sobre las rutas
finitas de \cref{x_reciprocidad:lib05:teo-rutas}, entonces
\begin{equation}
 \mathsf H(X\mid\Phi_n(X))=0.
 \label{x_reciprocidad:landauer:recuperacion-ruta}
\end{equation}
\end{corollary}
\begin{proof}
Escribamos la descomposición espectral finita
\(\rho=\sum_jp_j|v_j\rangle\langle v_j|\), con los \(v_j\)
ortonormales, \(p_j>0\) y \(\sum_jp_j=1\). La igualdad
\(Z^*Z=I\) implica que los \(Zv_j\) son ortonormales. Por tanto,
\(\rho_Z=\sum_jp_j|Zv_j\rangle\langle Zv_j|\) tiene los mismos
autovalores no nulos, lo que prueba la traza y la entropía.
Multiplicar por \(Z^*\) y \(Z\) recupera \(\rho\). Finalmente,
\(\Phi_n^{-1}\) recupera cada ruta de su pareja final; la
distribución condicional está concentrada en una sola ruta y su
entropía es cero.
\end{proof}

La conservación corresponde al archivo completo. Una lectura visible
\(q\), acompañada de memoria \(M\), recupera el estado bajo la
hipótesis de \eqref{x_reciprocidad:v:ant:eq:ii-kb-aut-fibra}; la información no
publicada por \(q\) permanece en
\(\mathsf H(M(X)\mid q(X))\). Descartar un complemento o un
terminal necesario define otra operación, cuya recuperabilidad se
determina sobre la información efectivamente retenida. La igualdad
espectral anterior conserva la distinción entre entropía del estado
completo, entropía de una reducción parcial y calor transferido.

En el modelo térmico de la sección
\ref{x_reciprocidad:v:ant:sec:ii-kb-aut-desarrollo}, con estados lógicos degenerados,
entorno isotermo y contabilidad de registros auxiliares, el término
asociado al borrado es
\(k_B\Theta\,\mathsf H(X\mid F(X))\), para un mapa lógico
determinista \(F\). Su contribución es nula para los transportes
recuperables anteriores; el reinicio de un TRIT uniforme satisface
\(Q_{\rm reset}\geq k_B\Theta\log3\). La preparación, el control,
la velocidad finita y el acoplamiento al entorno conservan sus
balances físicos propios. El calor total de una implementación
requiere componer esas contribuciones con la operación lógica.

\section{Conservación de observables y simetría relativa}
\label{x_reciprocidad:lib06:simetria}

El balance de la norma es el caso identidad de una ley operatoria
más informativa. Cuando los transportes son unitarios, con relativo
\(R\), el lector efectivo es
\begin{equation}
 \mathcal T_a(G)=\frac{a(a+1)G+R^*GR}{a(a+1)+1}.
 \label{x_reciprocidad:lib06:lector}
\end{equation}
El observable permanece fijo exactamente cuando \([G,R]=0\).
En el caso canónico,
\begin{equation}
 \mathcal T_8(G)-G=\frac{R^*GR-G}{73}.
 \label{x_reciprocidad:lib06:defecto-lector}
\end{equation}
Esta identidad determina tanto la conservación como la variación:
el defecto mide la diferencia entre dos realizaciones del mismo
lector, ponderada por el reparto \(72+1\).

La memoria completa transporta además el observable mediante
\(A\mapsto ZAZ^*\) en \(\operatorname{im}Z\). Los productos,
adjuntos y emparejamientos se conservan en esa imagen. La lectura
diagonal de canales y el observable transportado son operaciones
diferentes; el ejemplo del acarreo muestra la diferencia entre
\(74/73\) para la primera y el valor propio uno para el segundo.
La información permanece recuperable aunque una lectura diagonal
particular cambie de valor.

Las cargas variacionales de la \cref{x_reciprocidad:nuc11:noether} conservan su
fundamento en una acción y en un entrelazamiento explícitos. Para
\(A_{n+1}U_n=U_nA_n\), la acción de transporte da la carga
\(p_n^*A_nq_n\). Su conservación y la identidad de norma son
resultados compatibles, con hipótesis diferentes. De igual modo,
la conservación de una forma alternante o lorentziana en
\eqref{x_reciprocidad:lib04:balance-formas} mantiene su significado geométrico;
las estimaciones de ruido y contracción utilizan la positividad
de Hilbert declarada en sus pruebas.

\section{Función del registro en la transducción}
\label{x_reciprocidad:lib06:registro}

La constante holográfica de acoplamiento aritmético-geométrico tiene
una función definida por sus composiciones. El registro \(K\)
se obtiene previamente de la extracción orientada y de Hadamard.
Su lectura \(\kappa\) es reversible con base, longitud y origen
fijados. Su marco \(O_K\) identifica respuestas cíclicas porque
su Gram tiene una cota inferior estrictamente positiva. La
incidencia y la dirección geométrica utilizan sus lectores
declarados sobre el mismo registro.

La isometría de archivo conserva el Gram de ese marco, por lo que
mantiene las cotas de recuperación a cualquier profundidad. El
registro completo admite inversión de norma uno; las dos memorias
parciales conservan su factor \(9/4\) y su núcleo uniforme. La
carga adicional distingue la restitución del registro absoluto de
la de su componente centrada. La decodificación entera restituye
después la selección incidencial incluso en un umbral situado
exactamente sobre una coordenada.

Esta capacidad es más específica que conservar una magnitud total:
permite recuperar operaciones y relaciones que actúan sobre un
objeto generado. Sus garantías descansan en el Gram, en las
distancias de la imagen entera y en las inversas demostradas, no
en conservar por definición otra copia de la entrada. A la vez,
el número finito de ventanas de \(K\) permanece distinto de la
historia enriquecida completa; los otros registros y marcas
acompañantes mantienen sus propios transportes.

La órbita y sus transportes isométricos conservan la norma
\(\nu=\|K\|\). Normalizar sus columnas a norma uno exige
conservar \(\nu\) para invertir esa operación. La normalización
polar utiliza el par
\((F_K,G_K)\). La fórmula
\(O_K=F_KG_K^{1/2}\) mantiene orientación y amplitudes.
El lector posicional debe viajar con el cambio de carta, pues la
norma orbital común no vuelve invariante \(\kappa\) bajo una
renumeración aislada de sus componentes.

\section{Profundidad, dimensiones y realizaciones físicas}
\label{x_reciprocidad:lib06:profundidad}

La política bilateral que continúa un canal y archiva el otro
tiene agotamiento uniforme del terminal. Su límite es una
isometría hacia la suma de todas las memorias, y su inversa es
la serie adjunta convergente de
\eqref{x_reciprocidad:lib04:recuperacion-infinita}. Un número finito de
complementos posee dos reconstrucciones diferenciadas: el
adjunto truncado conserva un sesgo explícito y la inversa de
mínimos cuadrados lo elimina con la amplificación de ruido
calculada. Las evoluciones precedentes con terminal positivo
conservan esa componente en el archivo completo.

La dimensión del espacio de etiquetas, el rango de un proyector y
la dimensión de una realización física se relacionan mediante
sus mapas concretos. La unitaria de \(1001\) etiquetas y el
registro de doce ventanas son dos realizaciones explícitas de
dominios distintos. La conservación exterior del desarrollo
precedente prolonga los productos de Gram a cada grado finito,
incluyendo las contribuciones mixtas entre terminal y memorias.
Las realizaciones excepcionales y físicas permanecen vinculadas
a sus representaciones y operadores ya construidos; los rangos
numéricos por sí solos no sustituyen esa vinculación.

\Needspace{18\baselineskip}
\section{Conclusiones sobre la conservación evolutiva de la unidad}
Las acciones físicas y el protocolo gaussiano del manuscrito
conservan los acoplamientos, las formas, las unidades y las
preparaciones que intervienen en sus pruebas. El balance
bilateral aquí añadido amplía la familia de transportes
conservativos y sus lecturas; su aplicación a una realización
física se establece componiendo los operadores de esa
realización, sin reemplazarlos por una coincidencia de
normalización.

La unidad del desarrollo reside, por tanto, en una secuencia
demostrada: generación y completación, refinamiento con memoria,
realización unitaria, simetría relativa del lector y recuperación
de la estructura. La extrema y media razón holográfica recibe
así una exposición en la que cambio de escala, incidencia,
conservación y transducción poseen operaciones explícitas y
mantienen conjuntamente la información necesaria para invertir
sus lecturas.
